\section{Design e Implementazione}
È stata progettata un'architettura modulare e scalabile, con particolare attenzione alla gestione, al riuso della logica di gioco e della View.
Il package \emph{common} contiene il dominio condiviso, in modo che le due versioni \texttt{V1} e \texttt{V2} possano differire solo per il \textbf{Controller} e l'implementazione del \textbf{CollisionResolver}, difatto in fase di bootstrap, possono essere specificati, oltre alla strategia di generazione delle palline.

\subsection{Architettura del Sistema (MVC e Observer)}
L'implementazione per una visualizzazione reattiva \ref{sec:analisi_edt}, riprende \texttt{sketch02} e si basa su \textbf{Model-View-Controller (MVC)} con pattern \textbf{Observer}.

Per garantire la thread-safety, la rappresentazione visuale passa attraverso un \textbf{ViewModel}. Esso viene frequentemente aggiornato con dati differenti, rendendo cosi separati i dati dei model, soggetti ad update frequenti, e sincronizzazioni tra i vari thread.
In questo modo la View legge una vista coerente senza accedere direttamente alle strutture mutabili interne ai modelli.

\subsection{Gestione asincrona degli Eventi e Game Loop}
I \textbf{Controller} sono stati progettati per aggiornare in maniera costante la logica di gioco, gestire gli input e notificare la View.

Nella versione implementata con meccanismi low-level \texttt{V1}, usa \textbf{ActiveController}, il thread utilizza un \textbf{BoundedBuffer}, di dimensione limitata, per ricevere i comandi di gioco (input del giocatore).
Il controller non resta bloccato indefinitamente su una \emph{get()}, utilizza invece \emph{poll(waitMs)}.

Se un \textbf{GameCommand} è presente in coda, il comando viene eseguito; altrimenti il loop avanza sequenzialmente. In questo modo input e simulazione convivono nello stesso thread logico mantenendo una cadenza stabile.

La versione \texttt{V2} con \textbf{ExecutorController} mantiene un loop concettualmente simile, facendo però uso di costrutti high-level.
Ad esempio, al posto di \emph{BoundedBuffer} per i comandi, viene utilizzato \textbf{ConcurrentLinkedQueue}, usando \emph{offer} che a differenza della \emph{put} non è bloccante.

\subsection{Concorrenza delle Collisioni}
La risoluzione collisioni su $N$ palline è il punto in cui la concorrenza produce il massimo impatto, ma anche i principali rischi di contesa.
L'interfaccia \emph{resolveCollisions} è implementata in diverse versioni, di cui concorrenti: \textbf{MultithreadingCollisions} e \textbf{ExecutorCollisions}, mentre \textbf{SequentialCollisions} è la versione non concorrente, implementata in \texttt{sketch01}.

\subsubsection{Sincronizzazione tramite Barriera }
L'implementazione della classe \texttt{MultithreadingCollisions} implementata nella versione \texttt{V1}, utilizza meccanismi di low-level. I thread vengono istanziati una sola volta all'avvio e tenuti sempre in attesa tramite un ciclo continuo. 
Per coordinarli in modo sicuro con il game loop principale, sono state utilizzate due barriere cicliche.
La prima barriera funziona come un cancello di partenza, bloccando i worker finché il master non ha preparato e condiviso i dati del nuovo frame. Appena ricevono il via tramite la funzione \textbf{hitAndWaitAll()}, i thread si spartiscono in modo deterministico la lista delle palline elaborando le collisioni parallelamente, lavorando su porzioni separate senza ostacolarsi su strutture condivise.

Calcolate tutte le coppie di collisioni assegnate, ogni thread si ferma sulla seconda barriera. Solo quando tutti i partecipanti raggiungono questo traguardo finale, il master ha la certezza che la fisica del frame è completa e sblocca il ciclo per procedere al rendering. Questo design azzera l'enorme peso sul sistema operativo e sul Garbage Collector, permettendo al gioco di mantenere un frame rate fluido anche nella configurazione con migliaia di palline.

La \textbf{CyclicBarrier} è stata implementata sfruttando direttamente i monitor nativi di Java tramite i metodi \textbf{synchronized}, \textbf{wait()} e \textbf{notifyAll()}. La classe tiene traccia di quanti thread sono attesi, quanti sono effettivamente arrivati e a quale ciclo di esecuzione ci troviamo. Quando un worker arriva alla barriera, incrementa il contatore. Se mancano ancora altri thread all'appello, si mette in \textbf{wait()}, all'interno di un ciclo \textbf{while}, fondamentale per evitare che il thread riparta per sbaglio a causa di un risveglio spurio. Appena l'ultimo thread completa il suo lavoro, la barriera si resetta (definiamo cosi una \emph{barriera ciclica}), e chiama un \textbf{notifyAll()} per svegliare tutti i worker in attesa, facendoli ripartire.

\subsubsection{Ottimizzazione Locale}
La funzione \texttt{Ball.checkAndResolveCollisionSafe} esegue prima un pre-check geometrico senza lock. Questa scelta scarta il maggior numero di coppie, considerando solo quelle vicine. Riducendo così il numero di lock inutili su coppie sicuramente non in collisione.

\subsubsection{Lock Ordering: prevenzione da Deadlock}
Le collisioni deveno essere risolte garantendo la sezione critica su entrambe le palline, e più thread potrebbero tentare lock multipli in ordine inverso. Per evitare deadlock si applica un criterio deterministico di \textbf{Lock Ordering}.
Si confronta l'hash d'istanza della JVM fornito dal metodo:
\begin{verbatim}
if (System.identityHashCode(first) < System.identityHashCode(second))
\end{verbatim}
Ogni thread acquisisce prima il lock della pallina con identificativo inferiore e poi il secondo lock, in questo modo si elimina la possibilità di attese circolari.

\subsection{L'approccio Task-based, con meccanismi high-level}
L'evoluzione finale adotta un modello task-based utilizzando il Framework Java Executor. Vengono create micro-task e assegnate a un pool fisso di worker, la logica di suddivisione tra le coppie rimane la stessa della versione in \texttt{MultithreadingCollisions}. Ogni task elabora parallelamente e la sincronizzazione globale è ridotta alla sola barriera finale (\textbf{invokeAll}), tornando poi al flusso della Logica principale.

Questa scelta riduce sensibilmente la contesa sui monitor e sfrutta meglio i core disponibili, i meccanismi di high-level, riducono il rischio di errori di sincronizzazione, migliorano la manutenibilità del codice, e possiamo notare anche un miglioramento nella fluidità generale del gioco \ref{sec:benchmark_risultati}.